\subsection{The periodic law as a geometric coordinate}

The two-ended law has a further inverse consequence beyond normal
alternating channels.  Fix a clear nongrazing periodic polygon with $r$
distinct marked contacts.  Its collision points, tangent lines, phase
labels and closing translation are given; its contact curvatures and
normal graph germs are unknown.  At every physical phase retain the
conditional law of the two scaled tangent projections at two fixed,
distinct positive offsets.  These are function-valued data on a common
interior square, not samples of the unknown normal graph.

\begin{theorem}[Periodic contact reconstruction]\label{thm:v65-intro-inverse}
For the fixed marked polygon, the phase-resolved two-offset laws locally
determine every finite contact jet in the smooth category, including
curvature.  At an analytic base tuple they determine the complete contact
germs, and the future action map is a local biholomorphism in a
bounded-holomorphic norm on a small fixed disc.  Under a bounded local
analytic prior, real uniform density errors give conditional H\"older
bounds on a smaller disc.  If a fixed-lattice analytic table has every
labelled obstacle visited by the polygon, equality of these laws in the
local inverse class determines all of its obstacle images.
\end{theorem}

The separation of the two actions does not assume that they agree.
Writing the limiting density as
\[
 f_d(u,v)=Z_d^{-1}B^-(u)B^+(v)\{d-A(u)-C(v)\},
 \qquad A(0)=C(0)=0,
\]
the quotient $Q=f_{d_1}(u,v)f_{d_2}(0,0)/
[f_{d_2}(u,v)f_{d_1}(0,0)]$ gives
\[
 A(u)+C(v)=\frac{d_1d_2(1-Q)}{d_2-d_1Q}.
\]
Both flux factors and both normalizers cancel.  The oblique linear
momenta are then restored when passing from these physical actions to
the anchored variational actions.

There is also a geometric reason why all contact orders are invertible.
Let $\sigma_i$ be the signed one-step derivative of the future stable
half-line in normalized graph coordinates.  Dispersing positivity gives
$|\sigma_i|<1$.  At contact order $n$, the response is
\begin{equation}\label{eq:v65-intro-cayley}
 \mathcal C_n=(I+T_n)(I-T_n)^{-1},\qquad
 (T_nz)_i=\sigma_i^n z_{i+1}.
\end{equation}
This formula is obtained by counting the initial boundary visit once and
each internal visit twice in the actual stationary envelope.  It is not
an identification of a formal series with a smooth action.  The block
has determinant one for an even cycle; the explicitly computed odd-cycle
determinant is different but nonzero, including at odd contact orders.
The derivative of the full action is an invertible multiplier plus
contracted boundary evaluations.  A finite low-degree inverse and one
operator-norm tail inverse control all its lower-degree couplings.

\begin{proof}[Proof of Theorem~\ref{thm:v65-intro-inverse}]
Proposition~\ref{prop:v65-two-offset} separates the physical actions.
Lemma~\ref{lem:v65-envelope} proves smooth-jet factorization, and
Theorem~\ref{thm:v65-cyclic-jets} proves \eqref{eq:v65-intro-cayley}
and the local finite-jet inverse, beginning with its invertible
curvature block.  Lemma~\ref{lem:v65-holomorphic} constructs the
actual Banach-valued action map.  Its complete inverse, real-observation
estimate and exact analytic continuation consequence are respectively
Theorem~\ref{thm:v65-analytic-inverse},
Theorem~\ref{thm:v65-real-stability} and
Corollary~\ref{cor:v65-law-rigidity}.  All proofs are given in
Section~\ref{sec:v65-periodic-inverse}.  A nonnormal three-contact
family with independently varying curvatures is constructed in
Proposition~\ref{prop:v65-nonnormal-rigidity}.
\end{proof}

The relationship with variational and spectral methods is specific.
Hill's formula~\cite{Hill} concerns periodic Hessians and monodromy;
our forward normalization uses a Dirichlet corner cofactor before
conditioning.  The general-word length asymptotics and marked Lyapunov
recovery of B\'alint, De Simoi, Kaloshin and Leguil~\cite{BDKL}
concern periodic-orbit lengths, rather than endpoint-varying probability
laws.  Here the fixed marked polygon and the phase-resolved function data
make a separate contact inverse possible.  We do not identify these data
with an ordinary or enriched marked length spectrum.  In particular,
the enriched-spectrum rigidity theorem of Finamore and
Leguil~\cite{FinamoreLeguil} is a different observation model, not a
claimed consequence of the present theorem.  The matrix identity and
the functional-analytic inversion argument are not advanced as separate
novel rigidity theories; their geometric realization in the normalized
physical law and the actual contact action is the point of the argument.

The following alternating-channel results retain a different advantage:
they require only the gap and one law of each contact type, rather than
a fixed general polygon and two offsets per phase.  They also retain
the unknown-lattice reconstruction, finite global ambiguity and charged
finite-preparation results under their original hypotheses.  The new
periodic contact theorem extends the inverse domain without changing
those data requirements or replacing their conclusions.
